\documentclass[11pt, a4paper]{article}
\usepackage{amsmath,amssymb,amsfonts,amsthm}
\usepackage{geometry,hyperref,setspace,titlesec,booktabs,tabularx,array}
\geometry{left=2.5cm,right=2.5cm,top=2.5cm,bottom=2.5cm}
\setlength{\parindent}{0pt}\setlength{\parskip}{0.8em}\onehalfspacing
\hypersetup{colorlinks=true,linkcolor=blue,citecolor=blue,urlcolor=blue}
\newtheorem{proposition}{Proposition}
\title{\textbf{Language as an Activation-Control Signal for Associative Memory: \\ A Testable Framework from Propagation Terrain to Transformers}}
\author{Qing Zhou}\date{}
\begin{document}\maketitle


\begin{abstract}


This paper proposes a computational theory of language function: language is not only a symbolic sequence or an information carrier, but also a structured activation signal that controls associative-memory retrieval, alters internal-state propagation, and coordinates action. When two systems already share part of their experience, concepts, and behavioral rules, language need not transmit a complete world model item by item. A small number of cues can selectively activate existing memory routes in the receiver, allowing dispersed traces to converge, candidate interpretations to compete, and an actionable state to emerge.


Building on the view of memory as propagation terrain, the paper models linguistic input as an initial condition, boundary condition, and gating signal imposed on that terrain. It defines cue-entry match, retrieval gain, path ambiguity, convergence time, and behavioral effect, then combines them in a resource-constrained activation-cost functional. On the artificial-system side, modern Hopfield networks have a clear formal relation to attention. Contextual attention, feed-forward layers, residual streams, and layered transformation in Transformers can also be interpreted, respectively, as dynamic retrieval, parameterized long-term association, persistent state transport, and iterative constraint propagation. These correspondences, however, support only candidate mappings at the computational level. They do not by themselves prove that Transformers and brains are strictly isomorphic, or that the Transformer is the unique optimal architecture for language.


The paper's main contribution is to turn a strong intuition into a falsifiable framework. If language is an efficient activation-control signal for associative memory, then effective prompts should improve performance primarily by lowering the activation cost of correct routes, increasing the competitive advantage of target attractors, and reducing irrelevant spread. Multiple individually insufficient linguistic cues should sometimes produce superadditive retrieval. Shared background knowledge should reduce the linguistic length required for the same coordination task. Perturbing replay, contextual retrieval, or long-term associative structure should produce distinguishable deficits. The paper concludes with experimental designs for neuroscience, language-model analysis, and the training of knowledge-dense small models.


\end{abstract}
\textbf{Keywords:} language; associative memory; propagation terrain; activation control; modern Hopfield networks; Transformer; attention; prompting; shared knowledge; falsifiability



\section{The Problem: How Can a Small Signal Produce a Complex Internal Change?}


The sentence "the key is in the usual place" contains only a few words, yet it may be sufficient for a familiar listener to turn around, walk into a particular room, and open a particular drawer. For someone without the relevant shared history, the same sentence contains almost no executable content.


This example reveals a basic property of language: its effect depends not only on the signal, but also on the memory structure already present in the receiver. Language often does not transmit a complete model from scratch. It acts on shared knowledge, common experience, and present context to selectively activate what is already available.


Established traditions emphasize different aspects of language:


\begin{itemize}
\item Formal language theory studies symbolic composition and generative rules.
\item Information theory studies coding, noise, and transmission efficiency.
\item Statistical language modeling studies sequence probability and prediction.
\item Pragmatics studies speaker intention, common ground, and context.
\item Embodied cognition studies the relation of language to perception, action, and bodily state.
\end{itemize}


These traditions are not mutually exclusive. The present paper adds a middle-level computational question:


\begin{quote}
Once language enters a system with memory, goals, and the capacity to act, how does it change the propagation and convergence of internal activity?
\end{quote}


The central hypothesis is:


\begin{quote}
Language can be modeled as a structured activation signal that controls propagation through associative-memory terrain.
\end{quote}


Activation here does not mean merely increasing neural intensity, and control does not imply a central controller. More precisely, language changes which internal representations become accessible, which evidence is aggregated, which candidate states are suppressed, and which interpretation or action the system is more likely to enter.



\section{Theoretical Position: A Dynamics of Language Acting on Memory}


The paper does not reduce language to one mechanism. Language simultaneously has symbolic, statistical, social, perceptual, and action-oriented properties. Treating it as an activation-control signal addresses a narrower computational question: how language exploits a receiver's existing structure to produce an effect.


Three descriptive levels should be distinguished.


The first is the signal and symbolic level. Words, syntax, prosody, gesture, and writing provide recognizable linguistic structure. They determine how an input can be encoded, but they do not alone determine what the input will activate.


The second is the computational-mechanism level. Linguistic cues enter a system containing short-term state, long-term associations, context, goals, and gating. They alter retrieval, competition, and state transition. This is the main level addressed here.


The third is the implementation level. Brains may realize relevant functions through distributed activity, synaptic plasticity, working memory, hippocampal-cortical interaction, and control networks. Transformers realize some related computations through vector representations, attention, feed-forward layers, residual streams, and trained parameters. Similarity at the computational level does not imply identity of physical mechanism.


The framework therefore does not claim:


\begin{itemize}
\item that language is energy rather than information;
\item that all language understanding is memory retrieval;
\item that Attention corresponds item by item to biological synaptic activity;
\item that the Transformer has been proved to be the unique optimal architecture for language;
\item that the brain and the Transformer are globally related by a continuous invertible isomorphism;
\item or that formal similarity establishes common origin or mechanistic identity.
\end{itemize}


The proposal is a candidate model that can be compared with alternatives.



\section{From Transmitting Content to Activating Existing Structure}


\subsection{Shared Knowledge Determines Linguistic Compression}


Suppose a speaker wants a receiver to enter a target internal state $z^*$, such as making a judgment, recalling an event, or performing an action. The receiver already possesses a long-term memory structure $M$, the current context is $c$, and the linguistic signal is $\ell$.


Language understanding can be abstracted as


\[
z^* \approx \Psi(\ell,c;M),
\]

where $\Psi$ denotes the state evolution induced by language under a given memory structure and context.


If the receiver lacks the relevant prior structure $M$, a short linguistic signal will usually be unable to recover the target state. If speaker and receiver share a highly aligned structure, language needs only to distinguish among a small set of candidate routes. This motivates a conditional linguistic burden:


\[
L_{\min}(z^*|M,c,\epsilon)
=
\min_{\ell}
\left\{
|\ell|:
P\big(\Psi(\ell,c;M)=z^*\big)\ge 1-\epsilon
\right\}.
\tag{1}
\]

The quantity is the shortest linguistic length required to activate a target state with error no greater than $\epsilon$, given memory and context.


The model predicts that $L_{\min}$ decreases as shared knowledge becomes richer, the target route becomes more stable, and context becomes less ambiguous. When the internal structure is absent, language must carry a greater burden of teaching, definition, and disambiguation.


\subsection{A Linguistic Cue Is Not the Complete Memory}


A cue word may trigger the recollection of a complete event, but it does not contain all of that event's visual, emotional, and temporal detail. It functions more like a selective entry condition that moves the system from a broad space of possibilities into a relevant region.


Let $x$ be the current internal state and $\{z_j\}$ a set of candidate memory or action attractors. The entry match between linguistic cue $\ell$ and candidate state $z_j$ can be written as


\[
s_j(\ell,x,c)=K\big(e(\ell),k_j(x,c)\big),
\tag{2}
\]

where $e(\ell)$ is a linguistic encoding, $k_j$ is a state-dependent candidate-entry representation, and $K$ is a matching function.


The cue does not directly replace the system state with $z_j$. It changes the relative accessibility of candidate states:


\[
P(z_j|\ell,x,c,M)
\propto
\exp\left[
\beta s_j(\ell,x,c)
-R_j(M,x,c)
\right],
\tag{3}
\]

where $R_j$ is the effective propagation resistance of reaching $z_j$, and $\beta$ is sensitivity to entry match.


Language therefore derives its effect from both cue-entry match and the routes already shaped by history.


\subsection{Language Can Open, Close, and Reweight Routes}


Language does more than recall content. Negation, conditionals, role instructions, task rules, and metalinguistic prompts can reconfigure the interpretive frame.


For example, "treat it as a counterexample" may add no new fact, but it changes the relations onto which subsequent information is projected. This effect can be modeled as rapid gating of propagation-terrain parameters:


\[
M_{\text{eff}}=G(\ell,c,g)\odot M,
\tag{4}
\]

where $g$ denotes the current goal, $G$ is gating induced by language and context, and $\odot$ denotes selective modulation of connections, routes, or features.


Language therefore has at least three separable functions:


\begin{enumerate}
\item \textbf{Entry activation:} increasing access to selected memories or concepts.
\item \textbf{Constraint and disambiguation:} reducing candidate states inconsistent with context.
\item \textbf{Policy configuration:} changing the rules used in subsequent retrieval, reasoning, or action.
\end{enumerate}



\section{A Unified Dynamical Model of Linguistic Activation}


\subsection{State, Memory, and Linguistic Control}


Define a language-response system as


\[
\mathcal{LRS}=(\mathcal{X},M,E_\ell,G,\Psi,\Pi),
\]

where:


\begin{itemize}
\item $\mathcal{X}$ is the internal state space;
\item $M$ is propagation terrain formed by long-term experience;
\item $E_\ell$ encodes linguistic input as a control signal;
\item $G$ configures gating from language, goal, and context;
\item $\Psi$ is the state-propagation operator;
\item $\Pi$ maps internal state to observable output or action.
\end{itemize}


The dynamics are


\[
x_0=E_\ell(\ell,c),
\qquad
x_{t+1}
=
\Psi\big(x_t;M,G(\ell,c,g)\big)+\eta_t,
\tag{5}
\]

where $\eta_t$ represents noise, unmodeled input, or internal perturbation. The output is


\[
y=\Pi(x_T).
\tag{6}
\]

Language ability is expressed not only by next-token prediction, but also by whether linguistic input can move internal state into an appropriate region under finite time and resource constraints.


\subsection{Effective Activation Cost}


For a target state $z^*$, define the effective activation cost of a linguistic cue:


\[
\mathcal{A}(\ell,z^*|M,c)
=
\alpha_T T_{\text{conv}}
+\alpha_E E_{\text{act}}
+\alpha_S E_{\text{spread}}
+\alpha_W P_{\text{wrong}}
+\alpha_U H(Z|\ell,c,M).
\tag{7}
\]

Here:


\begin{itemize}
\item $T_{\text{conv}}$ is time to a stable output;
\item $E_{\text{act}}$ is total activation or computation cost;
\item $E_{\text{spread}}$ is irrelevant-route spread;
\item $P_{\text{wrong}}$ is the probability of convergence to an incorrect state;
\item $H(Z|\ell,c,M)$ is uncertainty over candidate states after the input.
\end{itemize}


An efficient linguistic signal should move the system into the target state at low cost without producing excessive incorrect activation.


\subsection{A Linguistic-Control Functional}


Under a language distribution $P_L$, task distribution $P_G$, and receiver-memory distribution $P_M$, define


\[
\mathcal{J}[E_\ell,G,\Psi,M]
=
\mathbb{E}_{\ell,g,M}
\left[
\mathcal{A}(\ell,z^*_g|M,c)
\right]
+\lambda_C\mathcal{C}
+\lambda_I\mathcal{I}
+\lambda_R\mathcal{R}
+\lambda_B\mathcal{B}.
\tag{8}
\]

The terms are:


\begin{itemize}
\item $\mathcal{C}$: structural and computational complexity;
\item $\mathcal{I}$: interference of new learning with old routes;
\item $\mathcal{R}$: fragility to paraphrase, noise, and context change;
\item $\mathcal{B}$: behavioral cost caused by false low-resistance routes, bias, or dangerous habits.
\end{itemize}


This functional does not assert that natural language has been proved globally optimal. It is a normative objective for comparing systems and training procedures under explicit constraints. Different systems may occupy different local solutions.


\subsection{Optimal Signals Depend on the Receiver}


The same sentence has different effects on different people, so an optimal linguistic signal cannot be defined independently of its receiver. For receiver $r$:


\[
\ell_r^*
=
\arg\min_{\ell\in\mathcal{L}}
\mathcal{A}(\ell,z^*|M_r,c_r).
\tag{9}
\]

Communication is therefore a receiver-model-dependent control problem. Differences in teaching, explanation, persuasion, and prompt design partly reflect how accurately the sender estimates the receiver's propagation terrain.



\section{Associative Memory, Attention, and Transformers}


\subsection{The Formal Relation Between Attention and Modern Hopfield Retrieval}


Classic Hopfield networks store patterns in an energy landscape and recover attractors from partial cues through iterative dynamics. Modern Hopfield networks extend storage and retrieval to continuous states and have a clear mathematical relation to softmax-based attention.


For query $q$, keys $K$, and values $V$, scaled dot-product attention is


\[
\operatorname{Attn}(q,K,V)
=
\operatorname{softmax}
\left(
\frac{qK^\top}{\sqrt d}
\right)V.
\tag{10}
\]

From an associative-memory perspective:


\begin{itemize}
\item $q$ is the current retrieval cue;
\item $K$ defines the addresses or match representations of candidate patterns;
\item softmax converts matches into competitive retrieval weights;
\item $V$ supplies the content retrieved and written into the current state.
\end{itemize}


This relation supports the computational interpretation that attention can perform associative retrieval. However, Transformer keys and values are jointly produced by present context and trained parameters; they are not identical to fixed memory patterns in a classic Hopfield network. An attention layer also need not iterate to an energy minimum.


\subsection{Contextual Attention as Dynamically Addressable Short-Term Association}


In self-attention, the current sequence provides a temporary associative store. Each token can access representations of other tokens according to its current query. This mechanism is well suited to reference resolution, dependency tracking, local composition, and contextual rule invocation.


If language is treated as an activation-control signal, self-attention corresponds to:


\begin{quote}
Linguistic cues dynamically selecting which currently active states should continue to influence subsequent computation.
\end{quote}


It primarily operates over memory visible in the present context rather than all long-term knowledge acquired during training.


\subsection{Feed-Forward Layers as Parameterized Association and Transformation}


A Transformer feed-forward layer is commonly written as


\[
\operatorname{FFN}(x)
=
W_2\sigma(W_1x+b_1)+b_2.
\tag{11}
\]

Prior work indicates that feed-forward layers can display key-value-memory-like behavior: input patterns activate intermediate units, which then contribute toward selected output directions. It remains too strong, however, to identify all FFN parameters directly with a long-term fact store. FFNs also perform nonlinear feature transformation, computation, and representational recoding.


A more conservative interpretation is:


\begin{quote}
The FFN is a trained structure for parameterized association and transformation, part of whose behavior can be operationalized as long-term associative memory.
\end{quote}


\subsection{Residual Streams as Persistent State Transport and Incremental Revision}


Residual connections allow each layer to add a correction to an existing state rather than reconstructing the entire state:


\[
h_{l+1}=h_l+\Delta_l(h_l).
\tag{12}
\]

From the propagation-terrain perspective, the residual stream preserves accessibility of the signal while retrieval, gating, and feature transformation incrementally revise a continuing state. This may reduce information loss and optimization difficulty in deep propagation, but it does not by itself establish a biological counterpart.


\subsection{Layered Computation as Iterative Constraint, Not Strict Monotonic Pruning}


As depth increases, a model can integrate broader context and recode candidate interpretations. The metaphor of progressively pruning a candidate space is useful, but real Transformer layers do not guarantee that the candidate set shrinks monotonically. Some layers may reintroduce, mix, or amplify information weakened earlier.


A more robust expression is


\[
Q_{l+1}
=
\mathcal{U}_l(Q_l,\ell,M_\theta),
\tag{13}
\]

where $Q_l$ denotes the candidate distribution implicit at layer $l$, and $\mathcal{U}_l$ may constrain, retrieve, reweight, or re-represent candidates. Successful reasoning requires the final distribution to support the output, not every layer to strictly reduce entropy.


\subsection{Component-Level Candidate Mapping}


\begin{center}\begin{tabularx}{\textwidth}{X X X}\toprule
\textbf{Function in the propagation-terrain framework} & \textbf{Candidate Transformer implementation} & \textbf{Status} \\
\midrule
Retrieval of a relevant state from a current cue & Query-Key-Value Attention & Clear formal relation \\
Temporary association within current context & Self-Attention context & Computational interpretation \\
Trained long-term association and transformation & FFNs and other parameters & Partial empirical support \\
Parallel retrieval in multiple relational subspaces & Multi-Head Attention & Candidate interpretation \\
Persistent transport and incremental revision of state & Residual Stream & Structural interpretation \\
Iterative integration and constraint propagation & Layer stacking & Candidate interpretation \\
Prompts changing retrieval and computation routes & In-context learning and prompting & Testable hypothesis \\
Offline experience resampling shaping long-term structure & Replay and continual learning during training & Engineering analogy \\
\bottomrule\end{tabularx}\end{center}


The mappings in this table have different strengths. Only the relation between attention and modern Hopfield retrieval is comparatively explicit. Other mappings require intervention experiments rather than conceptual analogy.



\section{Why Prompts Work: Activation, Disambiguation, and Program Configuration}


\subsection{Prompting Is Not Merely Adding Information}


Two prompts of similar length can produce very different performance. If the difference were determined only by new information content, semantically equivalent expressions should have approximately equivalent effects. In practice, systems can be highly sensitive to wording, example order, role instructions, and output format.


The propagation-terrain framework decomposes prompt effect as


\[
\Delta \operatorname{Perf}
=
\Delta A_{\text{entry}}
+\Delta A_{\text{disamb}}
+\Delta A_{\text{program}}
+\Delta A_{\text{verify}},
\tag{14}
\]

where:


\begin{itemize}
\item $\Delta A_{\text{entry}}$ improves access to relevant knowledge;
\item $\Delta A_{\text{disamb}}$ suppresses irrelevant interpretations;
\item $\Delta A_{\text{program}}$ configures intermediate computation or output policy;
\item $\Delta A_{\text{verify}}$ induces checking of candidate answers.
\end{itemize}


\subsection{Few-Shot Examples as Temporary Route Shaping}


Few-shot examples do more than display an answer. Within context, they establish local input-output relations that make subsequent queries more likely to propagate along the same structure. Their effect should depend on shared task structure rather than surface similarity alone.


The model predicts:


\begin{enumerate}
\item Structurally aligned examples with different vocabulary may outperform lexically similar examples with the wrong structure.
\item Mutually conflicting examples should increase route competition and output variance.
\item Example order should change temporary propagation terrain and therefore affect output.
\item When the necessary primitives are absent from long-term parameters, the benefit of few-shot prompting should be bounded.
\end{enumerate}


\subsection{Reasoning Prompts as Control of Computational Routes}


Instructions such as "check each condition," "first construct a counterexample," or "verify the units" may alter the organization of intermediate states. Such prompts need not supply new facts; they can guide the system onto a different computational route.


An effective reasoning prompt should lower failure probability along the correct route, but may increase computation length and activation cost. Its value should therefore be evaluated as


\[
\operatorname{Utility}(\ell)
=
\Delta P_{\text{correct}}
-\lambda_T\Delta T
-\lambda_E\Delta E.
\tag{15}
\]

A prompt that merely increases output length without changing internal route or correctness is not effective control.



\section{From Linguistic Communication to Shared Propagation Terrain}


\subsection{Common Ground Enables Highly Compressed Communication}


Communication efficiency depends strongly on shared concepts, rules, and experience. Define the effective shared terrain of speaker and receiver as


\[
\Omega(M_s,M_r,c)
=
\operatorname{Align}
\left(
\mathcal{R}(M_s,c),
\mathcal{R}(M_r,c)
\right),
\tag{16}
\]

where $\mathcal{R}(M,c)$ is the relational structure linguistically accessible in context.


When $\Omega$ is high, a short cue can reliably activate the target state. When it is low, the same cue is more likely to produce ambiguity, misunderstanding, or no useful response.


\subsection{Dialogue as Online Modeling of Another System's Terrain}


Speakers routinely adjust wording, examples, and depth in response to listener feedback. This can be written as


\[
\hat M_{r,t+1}
=
\operatorname{Update}
\left(
\hat M_{r,t},
y_t,
\ell_t
\right),
\tag{17}
\]

where $\hat M_{r,t}$ is the speaker's estimate of the receiver's memory structure and $y_t$ is receiver feedback.


Dialogue therefore not only exchanges content; it continually estimates and calibrates the activation entries shared by both parties. Explanatory ability can be understood as finding a low-cost, low-ambiguity activation route for a particular receiver.


\subsection{Metaphor and Analogy as Bridges Between Regions}


Metaphor and analogy map a familiar relational structure into an unfamiliar domain. When the mapping succeeds, it exploits an existing low-resistance route to establish new understanding quickly. When it is overextended, it imports incorrect constraints from the source domain.


Analogy quality should therefore be assessed by both structural preservation and false transfer:


\[
Q_{\text{analogy}}
=
\operatorname{Preserve}(\mathcal{S})
-\lambda_F\operatorname{FalseTransfer}(\mathcal{S}).
\tag{18}
\]


\section{Distinguishing Theoretical Propositions}


\subsection{Proposition One: Linguistic Efficiency Increases with Shared Knowledge}


Holding the target state and acceptable error fixed, the shortest linguistic length required for a task should decrease as effective shared structure increases:


\[
\frac{\partial L_{\min}}{\partial \Omega}<0.
\tag{19}
\]

The relation need not hold without limit. Shared bias or false prior structure may cause short language to activate an incorrect state rapidly.


\subsection{Proposition Two: Multiple Subthreshold Cues Can Produce Superadditive Retrieval}


If multiple cues each weakly activate a target route and jointly point toward the same state, their combination may cross a competitive threshold:


\[
P(z^*|\ell_1,\ell_2,M)
>
P(z^*|\ell_1,M)
+
P(z^*|\ell_2,M)
-P(z^*|M).
\tag{20}
\]

If cues point in conflicting directions, combination may instead reduce confidence or extend convergence time.


\subsection{Proposition Three: Effective Prompts Primarily Change Routes, Not Merely Token Count}


After prompt length and factual content are controlled, an effective prompt should alter intermediate representations, attention, or causal contribution routes, and lower the activation cost of the target state. If improvement is fully explained by added computation, the activation-control hypothesis adds no explanatory value.


\subsection{Proposition Four: Contextual-Retrieval and Long-Term-Association Damage Produce Different Deficits}


If contextual attention is selectively perturbed, a system should become worse at bindings and dependencies in the current sequence. If parameterized long-term associations are selectively perturbed, the system should lose trained knowledge or stable transformations while retaining some local contextual processing. A double dissociation would support the proposed functional division.


\subsection{Proposition Five: Language Learning Optimizes Accessibility, Not Only Frequency Fit}


When facts are held constant, training examples that provide more effective entries, structured contrasts, and cross-context bindings should allow a model to access knowledge with fewer examples and lower inference cost. The advantage should appear in prompt robustness, transfer, and interference resistance, not only training loss.



\section{Experimental Designs}


\subsection{Shared Knowledge and Minimum Linguistic Length}


Construct artificial agents with different background knowledge. All agents face the same target tasks, but their pretraining or memory stores contain different levels of shared structure. Ask a sender to generate the shortest instruction that allows a receiver to complete each task.


Measure:


\begin{itemize}
\item token count required for a fixed success rate;
\item instruction ambiguity;
\item convergence time or reasoning steps;
\item robustness to paraphrase and noise;
\item systematic misunderstanding caused by false shared priors.
\end{itemize}


If shared structure reduces instruction length beyond what word frequency explains, the result supports Proposition One.


\subsection{Subthreshold-Cue Aggregation}


Create target memories for which no single cue reliably retrieves the target, but two to four directionally aligned cues jointly identify it. Compare:


\begin{itemize}
\item single-cue retrieval;
\item simple concatenation of multiple cues;
\item a model with attention interactions removed;
\item top-k vector retrieval;
\item an iterative associative-retrieval model;
\item random or mutually conflicting cues.
\end{itemize}


The key criterion is whether combined benefit exceeds the linear sum of independent contributions and whether it is accompanied by measurable change in the competitive advantage of the target route.


\subsection{Causal-Route Experiments on Prompts}


Choose prompts with equivalent factual content but different control structure:


\begin{itemize}
\item answer directly;
\item list constraints before answering;
\item construct a counterexample before answering;
\item add irrelevant tokens only;
\item use a random paraphrase of equal length.
\end{itemize}


Use activation patching, attention intervention, causal tracing, layerwise probes, or intermediate-state decoding to determine whether improvement arises from repeatable internal-route changes. Attention visualization alone is insufficient to establish causality.


\subsection{Double Dissociation Between Attention and FFNs}


Apply controlled perturbations separately to attention outputs and FFN outputs, then evaluate two task classes:


\begin{enumerate}
\item tasks strongly dependent on current-context binding but requiring little long-term knowledge;
\item tasks with simple context but strong dependence on trained facts or transformations.
\end{enumerate}


If attention perturbation primarily harms the first class and FFN or parameterized-association perturbation primarily harms the second, the candidate functional division gains support. If damage patterns strongly overlap, the simple mapping should be revised.


\subsection{Accessibility-Oriented Training}


For small models of equal size, compare four training strategies:


\begin{enumerate}
\item random sample order;
\item ordering by knowledge dependency;
\item multiple-entry, multiple-context binding of the same knowledge;
\item multiple-entry binding with replay and interference monitoring.
\end{enumerate}


In addition to perplexity, evaluate:


\begin{itemize}
\item successful access under short prompts;
\item robustness to synonymous reformulation;
\item cross-task transfer;
\item old-knowledge retention after learning new material;
\item reasoning tokens and compute required for a correct answer.
\end{itemize}


If multiple-entry binding improves accessibility rather than simple repetition, its benefit should persist on unseen expressions and compositional tasks.


\subsection{Human Behavioral and Neural Experiments}


Human studies can manipulate common ground, cue combination, and goal-dependent gating while measuring response time, error type, confidence, and neural-representation change.


Key predictions include:


\begin{itemize}
\item shared background knowledge reduces linguistic length and response time;
\item multiple weak cues produce nonlinear retrieval improvement when presented together;
\item effective cues alter candidate competition before the target content is explicitly recalled;
\item the same word activates different relational structure under different goals;
\item learning and offline consolidation improve later cue efficiency, not only explicit-memory score.
\end{itemize}


These outcomes must be quantitatively compared with Bayesian inference, priming, working-memory, and general pragmatic models.



\section{Engineering Principles for Knowledge-Dense Small Models}


\subsection{Optimize Knowledge Entries, Not Only Knowledge Copies}


If language models often fail because knowledge exists but cannot be accessed reliably, training should optimize both storage and accessibility. For knowledge unit $k$, define entry coverage:


\[
\Gamma_k
=
\mathbb{E}_{\ell\sim P(\ell|k)}
\left[
P(\text{correctly access }k|\ell)
\right].
\tag{21}
\]

A knowledge-dense model should not merely retain knowledge in its parameters. It should allow a reasonable diversity of cues to access that knowledge.


\subsection{Multiple-Entry Binding}


The same knowledge should be bound through definitions, examples, counterexamples, causal relations, applications, and neighboring concepts. Multiple entries are not mechanical paraphrases. Each should provide a distinct but consistent relational constraint.


For each knowledge unit, an engineering system may maintain:


\begin{itemize}
\item a core proposition;
\item prerequisite concepts;
\item positive and negative examples;
\item common confusions;
\item executable applications;
\item boundaries with adjacent knowledge.
\end{itemize}


\subsection{Curriculum Ordering and Replay}


Stabilizing basic primitives before training combinations and cross-domain transfer reduces the risk that high-level knowledge rests on unstable entries. Replay should prioritize:


\begin{itemize}
\item old knowledge vulnerable to new-learning interference;
\item knowledge with low entry coverage;
\item knowledge with high dependency centrality across downstream tasks;
\item high-confidence but frequently incorrect knowledge;
\item samples with strong conflict between new and old routes.
\end{itemize}


\subsection{Distinguish Low Resistance from Correctness}


Training cannot merely make frequent routes easier to activate. False information, bias, and shortcut patterns can also become low-resistance routes. Counterexamples, source quality, calibration, conflict detection, and refusal mechanisms are therefore necessary.


Define route utility:


\[
U(p)
=
V_{\text{task}}(p)
-\lambda_F P_{\text{false}}(p)
-\lambda_H H_{\text{harm}}(p)
-\lambda_C C(p).
\tag{22}
\]

Low activation cost has engineering value only when route utility is positive.



\section{Falsification Conditions}


A theory has empirical content only if it can fail. The framework faces at least seven falsification conditions.


First, if common ground does not reduce the linguistic length or response time needed for a task after shared vocabulary, frequency, and task familiarity are controlled, the claim that language compresses control by exploiting existing terrain must be narrowed.


Second, if multiple individually subthreshold and directionally aligned cues always produce only linear benefit and never a stable competitive-threshold crossing, weak-cue field aggregation has no distinct value.


Third, if the full benefit of effective prompts is explained by additional tokens, compute time, or surface pattern matching, with no repeatable change in internal causal route, the interpretation of prompts as route configuration is unnecessary.


Fourth, if interventions on Attention and FFNs produce no distinguishable damage patterns, the interpretation of these components as dynamic retrieval and parameterized long-term association lacks support.


Fifth, if accessibility-oriented training improves only repetition of training expressions and does not improve paraphrase, compositional transfer, or interference resistance, multiple-entry binding is merely data duplication.


Sixth, if a simpler statistical language model or Bayesian pragmatic model explains all results better at equal complexity, the simpler account should be preferred.


Seventh, if any result can be explained by post-hoc adjustment of propagation terrain, resistance, or gating, the framework becomes unfalsifiable. Experiments must preregister operational variables, parameter ranges, and alternative models.



\section{Limitations and Boundaries}


This is a middle-level computational theory, not a complete theory of language or neuroscience. It does not explain the full generation of syntax, the social stabilization of meaning, language evolution, or the complete effects of emotional prosody, bodily action, and cultural institutions.


Propagation terrain, effective resistance, and activation cost are operational concepts, not independent physical quantities assumed to exist in the brain. They may be estimated from response time, activation threshold, transition probability, computation, and error spread, but no unique measurement currently exists.


The formal relation between Attention and modern Hopfield networks is real and important, but it does not imply that every Transformer computation is associative memory, or that biological memory uses the same algorithm.


Transformer success may jointly arise from training scale, optimization, data distribution, compatibility with parallel hardware, and architectural inductive bias. The statistical properties of language alone do not establish that Transformers are a unique or global optimum.


Language not only activates existing knowledge; it can teach new structure. An unfamiliar definition, logical proof, or precise instruction may create new internal relations during communication rather than merely retrieve old content. A complete theory requires both rapid structure formation and slow consolidation.


Finally, efficient activation does not guarantee truth, rationality, or benefit. Propaganda, bias, trauma cues, and false metaphors can all control internal state efficiently. Normative evaluation must remain independent of activation efficiency.



\section{Established Results, Model Definitions, and Hypotheses}


\begin{center}\begin{tabularx}{\textwidth}{X p{0.27\textwidth}}\toprule
\textbf{Claim} & \textbf{Status} \\
\midrule
Attention is a weighted value retrieval based on Query-Key match & Architectural fact \\
Modern Hopfield networks have a formal relation to softmax-based attention & Established theoretical result \\
FFNs can display key-value-memory-like behavior & Empirically supported interpretation \\
Activation cost consists of convergence time, spread, error, and uncertainty & Model definition \\
Language applies entry, constraint, and policy control to propagation terrain & Theoretical hypothesis \\
Shared knowledge reduces linguistic length for the same coordination task & Testable theoretical proposition \\
Directionally aligned subthreshold linguistic cues can produce superadditive retrieval & Testable theoretical proposition \\
Effective prompts improve performance by changing internal causal routes & Testable theoretical proposition \\
Attention primarily supports dynamic contextual retrieval, while FFNs partly support long-term association & Candidate functional division \\
The Transformer is the unique optimal implementation of linguistic propagation terrain & Not claimed \\
The brain and the Transformer are strictly dynamically isomorphic & Not claimed \\
Language only activates old knowledge and cannot create new structure & Explicitly rejected \\
\bottomrule\end{tabularx}\end{center}



\section{Final Thesis}


The framework can be condensed into the following loop:


\begin{quote}
Past experience shapes a system's associative structure and state-transition tendencies. Through symbols, order, context, and pragmatic constraint, language forms a structured control signal. That signal selectively activates some entries, suppresses some candidates, and configures subsequent computation. Internal activity competes and converges through the current propagation terrain, producing understanding, prediction, or action. The result of communication then changes future terrain through learning and replay.
\end{quote}


In one sentence:


\begin{quote}
The power of language lies not only in what it carries, but also in the state it can make a system with existing memory enter.
\end{quote}


This view offers a limited but testable bridge between human language research and Transformer analysis. The formal relation between Attention and modern Hopfield retrieval shows that associative retrieval is an important computational view of Transformers. Prompts, context, and layered transformation can then be studied as controls over internal propagation routes. Theoretical progress, however, does not come from declaring the two systems essentially identical. It comes from experiments that distinguish activation, retrieval, constraint, long-term association, and general statistical fitting.


If the framework is correct, building smaller and more reliable language models should not rely only on expanding parameters and corpora. It should also optimize knowledge entries, structured binding, curriculum order, offline replay, conflict detection, and causal auditability. A model's knowledge density depends not only on how much it stores, but also on whether it can access the right knowledge under the right context with low cost, low ambiguity, and low risk.



\section*{References}


Anderson, J. R. (1983). \textit{The Architecture of Cognition}. Harvard University Press.


Baddeley, A. (2003). Working memory: Looking back and looking forward. \textit{Nature Reviews Neuroscience}, 4, 829-839.


Buzsáki, G. (2015). Hippocampal sharp wave-ripple: A cognitive biomarker for episodic memory and planning. \textit{Hippocampus}, 25(10), 1073-1188.


Clark, A. (2013). Whatever next? Predictive brains, situated agents, and the future of cognitive science. \textit{Behavioral and Brain Sciences}, 36(3), 181-204.


Elhage, N., Nanda, N., Olsson, C., Henighan, T., Joseph, N., Mann, B., Askell, A., Bai, Y., Chen, A., Conerly, T., DasSarma, N., Drain, D., Ganguli, D., Hatfield-Dodds, Z., Hernandez, D., Jones, A., Kernion, J., Lovitt, L., Ndousse, K., et al. (2021). A mathematical framework for Transformer circuits. \textit{Transformer Circuits Thread}.


Friston, K. (2010). The free-energy principle: A unified brain theory? \textit{Nature Reviews Neuroscience}, 11, 127-138.


Geva, M., Schuster, R., Berant, J., \& Levy, O. (2021). Transformer feed-forward layers are key-value memories. \textit{Proceedings of EMNLP 2021}, 5484-5495.


Grice, H. P. (1975). Logic and conversation. In P. Cole \& J. L. Morgan (Eds.), \textit{Syntax and Semantics, Volume 3: Speech Acts}. Academic Press.


Hebb, D. O. (1949). \textit{The Organization of Behavior}. Wiley.


Hopfield, J. J. (1982). Neural networks and physical systems with emergent collective computational abilities. \textit{Proceedings of the National Academy of Sciences}, 79(8), 2554-2558.


McClelland, J. L., McNaughton, B. L., \& O'Reilly, R. C. (1995). Why there are complementary learning systems in the hippocampus and neocortex. \textit{Psychological Review}, 102(3), 419-457.


Petroni, F., Rocktäschel, T., Riedel, S., Lewis, P., Bakhtin, A., Wu, Y., \& Miller, A. H. (2019). Language models as knowledge bases? \textit{Proceedings of EMNLP-IJCNLP 2019}, 2463-2473.


Ramsauer, H., Schäfl, B., Lehner, J., Seidl, P., Widrich, M., Adler, T., Gruber, L., Holzleitner, M., Pavlović, M., Sandve, G. K., Greiff, V., Kreil, D., Kopp, M., Klambauer, G., Brandstetter, J., \& Hochreiter, S. (2021). Hopfield networks is all you need. \textit{International Conference on Learning Representations}.


Sperber, D., \& Wilson, D. (1995). \textit{Relevance: Communication and Cognition} (2nd ed.). Blackwell.


Vaswani, A., Shazeer, N., Parmar, N., Uszkoreit, J., Jones, L., Gomez, A. N., Kaiser, Ł., \& Polosukhin, I. (2017). Attention is all you need. \textit{Advances in Neural Information Processing Systems}, 30.


Wilson, M. A., \& McNaughton, B. L. (1994). Reactivation of hippocampal ensemble memories during sleep. \textit{Science}, 265(5172), 676-679.



\end{document}